\clearpage
\section{047 -- Stackup JLC08161H-2116 e piste RF microstrip \SI{50}{\ohm} a \SI{0,20}{\milli\metre}}

\begin{tabularx}{\linewidth}{@{}>{\bfseries}l X@{}}
\toprule
Data & 11/11/2026\\
Revisione & non indicata\\
Redatto da & Claude Code\\
Approvazione & Brik81\\
Stato & \tg{tagdone}{MODIFICA ESEGUITA} \quad \tg{tagwarn}{CRITICITÀ APERTE}\\
\bottomrule
\end{tabularx}

\subsection{Componenti, net e fogli coinvolti}
\begin{itemize}
  \item File: \code{SYMPHONY.kicad\_pcb} (stackup), \code{SYMPHONY.kicad\_dru} (regole). \code{SYMPHONY.kicad\_pro} \emph{non} modificato.
  \item Net class \code{RF\_50R\_CPWG}: \code{RF\_ANT}, \code{RF\_SW\_*}, \code{RF\_RX\_F1}\ldots\code{F4}, \code{RF\_LNA\_*}, \code{RF\_RX\_MATCH}, \code{RF\_SI\_TX}, \code{RF\_TX\_M1}/\code{M2}, \code{RF\_ATT\_*}, \code{RF\_PA\_IN}, \code{RF\_PA\_MATCH}, \code{RF\_LPF\_*}.
  \item Net class \code{RF\_PA\_OUT} (larghezza \SI{1,0}{\milli\metre}), net \code{/UHF/RF\_ANT}.
  \item Layer: F.Cu (piste RF), In1 (riferimento GND; zone GND anche su In3 e In6).
  \item Backup prima della modifica: \code{report\_src/bak\_stackup\_2116/} (\code{.kicad\_pcb}, \code{.kicad\_dru}, \code{.kicad\_pro}).
\end{itemize}

\subsection{Cosa è cambiato}
\textbf{Contesto.} Le piste RF sono su F.Cu senza massa coplanare sul top (nessuna zona GND su F.Cu): sono quindi \emph{microstrip} riferite a In1. Brik81 vuole \SI{50}{\ohm} esatti e larghezza fissa \SI{0,20}{\milli\metre} (minimo ECSS esterni nel \code{.kicad\_dru}). Il progetto per ora \textbf{non prevede soldermask} (decisione da confermare con i responsabili).

\textbf{Calcolo} (Hammerstad-Jensen con correzione dello spessore del rame \SI{35}{\micro\metre}, senza soldermask), pista \SI{0,20}{\milli\metre} sugli stackup JLC a 8 layer da \SI{1,6}{\milli\metre}:
\begin{center}
\begin{tabular}{@{}lccc@{}}
\toprule
Stackup JLC & $h$ [\si{\milli\metre}] & $\varepsilon_r$ & $Z_0$\\
\midrule
1080 & \num{0,0764} & \num{3,91} & \SI{40}{\ohm}\\
\textbf{2116 (scelto)} & \num{0,1164} & \num{4,16} & \textbf{\SI{50,8}{\ohm}}\\
1080$\times$2 & \num{0,1604} & \num{3,91} & \SI{62}{\ohm}\\
7628 & \num{0,2104} & \num{4,4} & \SI{68}{\ohm}\\
\bottomrule
\end{tabular}
\end{center}
Con soldermask JLC ($\varepsilon_r=3{,}8$) l'impedenza scende di circa \SI{1}{\ohm}. Fonte stackup: repository gsuberland/jlcpcb\_autogenerated\_stackups (KiCad); $\varepsilon_r$ da jlcpcb.com/impedance.

\textbf{1) Stackup in \code{SYMPHONY.kicad\_pcb}} (Board Setup, Physical Stackup). Spessori in \si{\milli\metre}.
\begin{center}
\small
\begin{tabular}{@{}p{2.6cm}p{6.2cm}p{6.0cm}@{}}
\toprule
Elemento & Prima (a mano, nessuno stackup JLC) & Dopo (JLC08161H-2116)\\
\midrule
F.Cu & \num{0,035} & \num{0,035}\\
Dielettrico 1 & prepreg 3313 \num{0,1} + sottostrato 3313 \num{0,0994} (circa \num{0,2}, $\varepsilon_r$ \num{4,1}) & PP 2116 \num{0,1164}, $\varepsilon_r$ \num{4,16}\\
In1 & \num{0,0152} & \num{0,0152}\\
Dielettrico 2 & core \num{0,2438}, $\varepsilon_r$ \num{4,6} & core \num{0,3}, $\varepsilon_r$ \num{4,41}\\
In2 & \num{0,0152} & \num{0,0152}\\
Dielettrico 3 & prepreg 3313 \num{0,1} + \num{0,107} & PP 1080$\times$2 \num{0,1528}, $\varepsilon_r$ \num{3,91}\\
In3\ldots In6 & \num{0,035} ciascuno & \num{0,0152} ciascuno\\
Dielettrici 4\ldots 7 & core \num{0,2438} ($\varepsilon_r$ \num{4,5}) / prepreg \num{0,1} ($\varepsilon_r$ \num{4,5}) / core \num{0,2438} / prepreg \num{0,1} ($\varepsilon_r$ \num{4,5}) & core \num{0,3} / PP 1080$\times$2 \num{0,1528} / core \num{0,3} / PP 2116 \num{0,1164}, con $\varepsilon_r$ \num{4,41} (core), \num{3,91} (1080$\times$2), \num{4,16} (2116)\\
B.Cu & \num{0,035} & \num{0,035}\\
Spessore generale & \num{1,6082} & \num{1,6}\\
\bottomrule
\end{tabular}
\end{center}
Sequenza finale: F.Cu / PP 2116 / In1 / core / In2 / PP 1080$\times$2 / In3 / core / In4 / PP 1080$\times$2 / In5 / core / In6 / PP 2116 / B.Cu. Maschere lasciate nello stackup (non influiscono). \code{dielectric\_constraints} (impedance controlled) già attivo. Finitura ENIG invariata.

\textbf{2) Regole in \code{SYMPHONY.kicad\_dru}.}
\begin{center}
\small
\begin{tabular}{@{}p{3.4cm}p{5.8cm}p{5.8cm}@{}}
\toprule
Voce & Prima & Dopo\\
\midrule
Intestazione & 4 layer JLC04161H-3313A & 8 layer \SI{1,6}{\milli\metre} JLC08161H-2116\\
Sezione 7 RF & regola ``RF CPWG gap 0.25'' (CPWG w \SI{0,32}{\milli\metre} / gap \SI{0,25}{\milli\metre}, valori del vecchio stackup) & rimossa\\
Nuova regola & non presente & ``RF microstrip larghezza 0.20'': layer esterni, piste \code{RF\_50R\_CPWG}, \code{track\_width} min \SI{0,19}{\milli\metre} / opt \SI{0,20}{\milli\metre} / max \SI{0,21}{\milli\metre}\\
Nuova regola & non presente & ``RF microstrip clearance 0.35'': piste \code{RF\_50R\_CPWG} verso altre net $\ge\SI{0,35}{\milli\metre}$ ($\approx 3h$, per non diventare coplanare)\\
Regola \code{RF\_50R\_MS} (classe non assegnata) & keep-out \SI{0,65}{\milli\metre} & keep-out \SI{0,35}{\milli\metre}\\
\bottomrule
\end{tabular}
\end{center}
Il nome della net class resta \code{RF\_50R\_CPWG} per non rompere le assegnazioni, ma la struttura è una microstrip.

\textbf{3) \code{SYMPHONY.kicad\_pro} non modificato} (Schematic Editor ancora aperto, rischio di sovrascrittura): la net class \code{RF\_50R\_CPWG} ha già \code{track\_width} \SI{0,2}{\milli\metre} e clearance di classe \SI{0,25}{\milli\metre}; la regola custom da \SI{0,35}{\milli\metre} prevale.

\textbf{Verifica.} DRC con kicad-cli prima/dopo: nessuna nuova violazione, tranne 2 \code{track\_width} su \code{/UHF/RF\_ANT}: due spezzoni F.Cu da \SI{0,5}{\milli\metre}, lunghi \SI{0,007}{\milli\metre} e \SI{0,113}{\milli\metre}, vicino a x \num{114,55} y \num{82,42}, residui al pad del connettore. Le altre circa 567 violazioni e 423 non collegati esistevano già (routing in corso).

\subsection{Motivazione}
Larghezza fissa \SI{0,20}{\milli\metre} e \SI{50}{\ohm} richiesti da Brik81; con pista \SI{0,20}{\milli\metre} lo stackup 2116 (\SI{50,8}{\ohm}) è quello, tra i JLC a 8 layer da \SI{1,6}{\milli\metre} calcolati, più vicino a \SI{50}{\ohm}.

\textbf{Approvazione di Brik81 (testuale):} ``fissiamo la traccia a 0.2 ohm [mm], che modo abbiamo per tunare la ipedneza a 50?''; e poi, dopo la proposta step-by-step dello stackup 2116: ``ascolta ho choiuso ilk pcb non ho tmepo puoi fare te perfavore?''.

Le modifiche sono state eseguite dall'agente principale.

\subsection{Impatti e criticità}
\begin{critbox}
\textbf{RF\_PA\_OUT (corretta nella entry 048).} L'ipotesi iniziale di circa \SIrange{20}{22}{\ohm} da rivedere era un errore di analisi: \code{RF\_PA\_OUT} è il nodo di uscita del PA prima della rete di adattamento, a bassa impedenza, e la larghezza \SI{1,0}{\milli\metre} è corretta. Nessuna modifica.
\end{critbox}
\begin{notebox}
\textbf{RF\_ANT.} I due spezzoni da \SI{0,5}{\milli\metre} vanno sistemati.\\
\textbf{Massa.} In1 GND deve restare pieno sotto tutte le piste RF (e In6 se ci sono RF sul bottom).\\
\textbf{Soldermask.} Sì/no da decidere con i responsabili. Senza coating il \code{.kicad\_dru} richiede clearance esterna \SI{0,25}{\milli\metre} (regola NO COATING, attualmente commentata).\\
\textbf{Stima precedente superata.} La stima di \SI{0,17}{\milli\metre} data in precedenza a Brik81 era basata su una lettura errata dello stackup ($h=\SI{0,1}{\milli\metre}$ invece di circa \SI{0,2}{\milli\metre} con il sottostrato).
\end{notebox}

\subsection{Verifiche da fare}
\begin{itemize}
  \item Verificare l'impedenza sul calcolatore JLC.
  \item Ordinare con Impedance control = Yes, stackup JLC08161H-2116, 8 layer, \SI{1,6}{\milli\metre}, FR-4 TG155, rame 1 oz esterno / 0,5 oz interno, nota: ``50 ohm single-ended microstrip, track 0.20 mm, L1 ref L2, L8 ref L7''.
  \item Sistemare i due spezzoni da \SI{0,5}{\milli\metre} di \code{RF\_ANT}.
  \item (Chiusa nella entry 048: larghezza di \code{RF\_PA\_OUT} confermata a \SI{1,0}{\milli\metre}.)
  \item Decidere con i responsabili soldermask sì/no.
\end{itemize}
